\section{Radius-uniform cross-curve distortion on every real word}
\label{sec:xd-cross-curve}

The transverse comparison required for a whole-cylinder pressure sum
is stronger than distortion along one short curve.  We prove that
comparison here directly from the already established variational
localization of module 308.  In particular, the present argument does
not assume that a physical word domain is convex, that two points can
be joined through regular physical trajectories, or that a homogeneous
planar cell has any unproved two-sided diameter property.
The new conclusion is a uniform comparison of actual arclength
Jacobians at different physical points and different cone directions.
The localization input and its anchor-layer proof are inherited.

All trajectories below belong to the original unit-speed triangular
Lorentz gas, with
\[
 R_-={9\over20},\quad R_+={47\over100},\quad
 \tau_-={3\over50},\quad
 a={100\over47},\quad K={170\over9},\quad
 \varkappa_+={20\over9},\quad \vartheta_s={47\over48}.
\]
Let $\tau_+<\infty$ be the uniform physical flight bound proved in
module 282.  As in module 313, tangent norms are the Euclidean norms
in $(s,\varphi)$, where $s=R\alpha$ and $p=\sin\varphi$.
There is no measure normalization in a tangent Jacobian.
The homogeneity partition is the original $q=2$ partition with one
fixed $k_0\ge2$:
\[
 H_k^+=\{\pi/2-k^{-2}<\varphi<\pi/2-(k+1)^{-2}\},\qquad
 H_k^-=-H_k^+\quad(k\ge k_0),
\]
and $H_0$ is the intervening central strip.  Including strip boundaries
by a fixed disjoint convention changes none of the inequalities below;
true grazing points, at which the requested derivative is undefined,
are excluded.

Put
\begin{equation}\label{eq:xd-constants-first}
 N=94,\qquad q_{\rm a}={3263680\over3263681},\qquad
 q=q_{\rm a}^{1/(4N)},\qquad C=10^8,
\end{equation}
where $q_{\rm a},N,C$ are exactly the constants of module 308.
In particular $q$ is not substituted for the original guard constant
$q_0=47/53$.  Define the following finite, collision-independent
constants:
\begin{align}
 C_Q&=2(2R_+C)^{1/4}q_{\rm a}^{-1/2},
 &C_\tau&=(1+q^{-1})C_Q,\label{eq:xd-CQ}\\
 C_\varphi&={\pi\over2}
       \left({C_Q\over R_-}+{2C_\tau\over\tau_-}\right),
 &D_c&={2C_\varphi\over1-q},\qquad
 D_\tau={2C_\tau\over1-q},\label{eq:xd-Cphi}\\
 D_v&={K-a+
       (\tau_-^{-1}+(2a\tau_-^2)^{-1})D_c
                    +\tau_-^{-2}D_\tau\over1-\vartheta_s},
 \label{eq:xd-Dv}\\
 H_{\log}&=\max\left\{2\pi3^{2/3},
                         {\pi^{5/3}\over2}k_0^2\right\},
 &D_{\log}&={2H_{\log}C_\varphi^{1/3}\over1-q^{1/3}},
 \label{eq:xd-Dlog}\\
 D_A&={ (\varkappa_++K)D_\tau+\tau_+D_v+D_c
                                                 \over2a\tau_-},
 &D_{\rm end}&=\log{1+K^2\over1+a^2},\label{eq:xd-DA}\\
 D_\times&=D_A+D_{\log}+D_{\rm end},
 &C_\times&=\exp(D_\times).
 \label{eq:xd-final-constant}
\end{align}
These conservative constants may be very large.  Their finiteness is
uniform in the radius and every physical collision count; no assertion
of a useful finite-sample numerical size is needed.

\subsection{A two-record consequence of the inherited anchor theorem}

\begin{lemma}[Paired real records, including disconnected word components]
\label{lem:xd-paired-records}
Let $m\ge1$ and let $Q,\widehat Q$ be any two regular physical
$m$-flight records at the same radius and with the same ordered lifted
center word.  Choose the common lift $O_0=0$ and write
\[
 b_i=q^i+q^{m-i}\qquad(0\le i\le m).
\]
For the actual contacts, flight lengths, incidences, and collision
angles one has
\begin{align}
 |Q_i-\widehat Q_i|&\le C_Qb_i,
                                      &&0\le i\le m,
                                  \label{eq:xd-contact}\\
 |\tau_i-\widehat\tau_i|&\le C_\tau b_i,
                                      &&0\le i<m,
                                  \label{eq:xd-flight}\\
 |\varphi_i-\widehat\varphi_i|
       &\le C_\varphi b_i,
                                      &&0\le i\le m,
                                  \label{eq:xd-phase}\\
 \sum_{i=0}^m|c_i-\widehat c_i|&\le D_c,
 \qquad
 \sum_{i=0}^{m-1}|\tau_i-\widehat\tau_i|\le D_\tau.
                                  \label{eq:xd-summed-data}
\end{align}
In particular the actual two-sided phase distance is exponentially
small toward the middle of a complete common physical word.
The assertion requires no homogeneity or endpoint cap, and no minimum
incidence or clearance.  It compares two real records; the auxiliary
minimum used in its proof is not a replacement trajectory.
\end{lemma}

\begin{proof}
Fix any one global minimum $Y$ of the disk-product problem for this
word and radius.  It is a common comparison for both records; the
proof does not require uniqueness of that minimum.  In the notation of
Corollary~\ref{cor:ae-geometric}, let
$d_i=\min(i,m-i)$ and
$j(d)=\max\{0,\lfloor d/N\rfloor-1\}$.
The already proved quartic position budget, applied separately to the
two real references and to this same $Y$, gives
\[
 |Q_i-Y_i|^4\le C E_R(Q)q_{\rm a}^{j(d_i)},\qquad
 |\widehat Q_i-Y_i|^4
                 \le C E_R(\widehat Q)q_{\rm a}^{j(d_i)}.
\]
Both original endpoint deficits are at most $2R\le2R_+$, including
when $i$ is an endpoint.  Since $j(d)\ge d/N-2$,
\[
 |Q_i-\widehat Q_i|
 \le2(2R_+C)^{1/4}q_{\rm a}^{j(d_i)/4}
 \le C_Q q^{d_i}\le C_Q b_i.
\]
This is a consequence of the full inherited positive energy estimate,
not an integration of a local endpoint derivative across a singular
boundary.  It proves \eqref{eq:xd-contact}.

The triangle inequality gives
\[
 |\tau_i-\widehat\tau_i|
 \le|Q_i-\widehat Q_i|+|Q_{i+1}-\widehat Q_{i+1}|.
\]
Use $b_{i+1}\le q^{-1}b_i$ to prove \eqref{eq:xd-flight}.
For nonzero vectors of lengths at least $\tau_-$, their normalized
directions differ by at most $2/\tau_-$ times their difference.
Thus the unit direction of either flight adjacent to contact $i$
differs from its paired direction by at most
$2C_\tau b_i/\tau_-$; for the preceding flight use
$b_{i-1}\le q^{-1}b_i$ directly in the sum of its endpoint bounds.
The contact normals differ by at most $C_Qb_i/R_-$.

At $i<m$, the unit vector $(c_i,p_i)$ is obtained by resolving the
actual outgoing velocity in the normal and tangent frame of $Q_i$.
The difference of the two resolved unit vectors is consequently at
most the sum of the velocity and frame differences.  At $i=m$, use
the incoming resolution $(-v_{m-1}\cdot n_m,v_{m-1}\cdot t_m)$;
it gives the same bound.  Both resolved vectors lie on the right
closed semicircle, with angles in $[-\pi/2,\pi/2]$.  The angular
difference on that semicircle is at most $\pi/2$ times their Euclidean
chord distance.  This proves \eqref{eq:xd-phase}, without an
inverse-incidence differentiation of $\arcsin p$.
Likewise the shorter arclength displacement of the contact on its
radius-$R$ circle is at most $\pi|Q_i-\widehat Q_i|/2$, so the same
argument controls the actual phase metric modulo the circle seam.

Finally, $|c_i-\widehat c_i|\le|\varphi_i-\widehat\varphi_i|$ and
$\sum_{i=0}^m b_i\le2/(1-q)$.  Summing the established inequalities
proves \eqref{eq:xd-summed-data}.  No connecting physical path or
two-dimensional cell shape has appeared in the argument.
\end{proof}

\begin{lemma}[The logarithmic incidence estimate on the actual strips]
\label{lem:xd-log-incidence}
For two angles in the closure of the same member of the fixed $q=2$
homogeneity partition,
\begin{equation}\label{eq:xd-log-incidence}
 |\log\cos\varphi-\log\cos\widehat\varphi|
                \le H_{\log}|\varphi-\widehat\varphi|^{1/3}.
\end{equation}
For any subset $J\subset\{0,\ldots,m\}$ of matching homogeneous
indices of the records in Lemma~\ref{lem:xd-paired-records},
\begin{equation}\label{eq:xd-log-sum}
 \sum_{i\in J}|\log c_i-\log\widehat c_i|\le D_{\log}.
\end{equation}
\end{lemma}

\begin{proof}
In $\overline H_k^\pm$ put $\delta=|\varphi-\widehat\varphi|$.
Then $\delta\le3k^{-3}$ and
$\cos\varphi,\cos\widehat\varphi\ge(2/\pi)(k+1)^{-2}$.
The mean value bound for $\log\cos\varphi$ gives
\[
 |\log c-\log\widehat c|
 \le{\pi\over2}(k+1)^2\delta
 \le2\pi k^2\delta
 \le2\pi3^{2/3}\delta^{1/3}.
\]
For the central strip $c\ge\sin(k_0^{-2})\ge(2/\pi)k_0^{-2}$
and $\delta\le\pi$.  Hence the same mean value argument gives the
other constant in \eqref{eq:xd-Dlog}.  Both calculations hold on the
closed strips.  The original disjoint assignment of their boundary
points is retained; no symmetry of that assignment is required.
For paired records use \eqref{eq:xd-phase},
$(A+B)^{1/3}\le A^{1/3}+B^{1/3}$, and the two geometric sums:
\[
 \sum_{i\in J}|\log c_i-\log\widehat c_i|
 \le H_{\log}C_\varphi^{1/3}
        \sum_{i=0}^m(q^{i/3}+q^{(m-i)/3})
 \le {2H_{\log}C_\varphi^{1/3}\over1-q^{1/3}}.
\]
The common strip is used only where its logarithm is actually summed.
\end{proof}

\subsection{Cross-direction distortion and the exact endpoint factor}

For a nonzero vector $\xi$ at $x$ of slope in $[a,K]$, write
\[
 \mathcal J_m(x,\xi)={\|D T_R^m(x)\xi\|\over\|\xi\|}.
\]
This is an actual forward arclength expansion in the unstable cone.
When $\xi$ is tangent to a smooth curve at $x$, it is precisely the
Jacobian along that curve at that point.  Curvature and a relation
between two such curves are not needed for this pointwise comparison.

\begin{theorem}[Uniform distortion across all physical components]
\label{thm:xd-cross-distortion}
Let $x,\widehat x$ generate two regular $m$-flight records with the
same radius and the same full lifted physical word, and let
$\xi,\widehat\xi$ be arbitrary nonzero vectors with slopes in
$[a,K]$ at the two initial points.  Suppose that the actual collision
angles at each interior index $1\le i<m$ lie in the closure of the
same respective homogeneity strip.  There is then the endpoint-exact bound
\begin{equation}\label{eq:xd-weighted-distortion}
 \left|\log{c_m(x)\mathcal J_m(x,\xi)
              \over c_m(\widehat x)\mathcal J_m(\widehat x,\widehat\xi)}
                                                        \right|
 \le D_\times.
\end{equation}
Equivalently,
\begin{equation}\label{eq:xd-endpoint-residual}
 e^{-D_\times}{c_m(\widehat x)\over c_m(x)}
 \le {\mathcal J_m(x,\xi)\over
             \mathcal J_m(\widehat x,\widehat\xi)}
 \le e^{D_\times}{c_m(\widehat x)\over c_m(x)}.
\end{equation}
No homogeneity assumption is made on either endpoint in these two
assertions.  If the terminal angles also lie in one common closed strip,
the same displayed constant gives the unweighted bound
\begin{equation}\label{eq:xd-unweighted-distortion}
 \left|\log{\mathcal J_m(x,\xi)
                   \over\mathcal J_m(\widehat x,\widehat\xi)}\right|
                                                         \le D_\times.
\end{equation}
The constants are uniform in the radius, collision count, complete
physical word, and the two cone directions.  In particular the two
points may belong to different connected regular components of that
word, including components separated by true unvisited tangencies
or any of the unchanged section and exact-label tests.
\end{theorem}

\begin{proof}
Let $v_i,\widehat v_i\in[a,K]$ be the slopes of the successive
actual images of the two vectors, and put $\varkappa=1/R$.
The exact signed differential in Lemma~\ref{lem:eo-differential}
has a scalar minus sign; that sign changes neither slope nor
Euclidean length.  Its positive scalar numerator and slope map are
\begin{equation}\label{eq:xd-A-and-f}
 A_i=\tau_i(\varkappa+v_i)+c_i,\qquad
 v_{i+1}=f(v_i;c_i,c_{i+1},\tau_i),\qquad
 f(v;c,d,\tau)=\varkappa+
                     {d(\varkappa+v)\over\tau(\varkappa+v)+c}.
\end{equation}
Here $A_i\ge2a\tau_->0$.  The native contraction calculation gives
$0\le\partial_v f\le\vartheta_s=47/48$.
On the full parameter rectangle
\[
 v\in[a,K],\quad c,d\in[0,1],\quad\tau\in[\tau_-,\tau_+]
\]
the other exact derivatives satisfy
\begin{equation}\label{eq:xd-parameter-derivatives}
 |\partial_c f|\le{1\over2a\tau_-^2},\qquad
 |\partial_d f|\le{1\over\tau_-},\qquad
 |\partial_\tau f|\le{1\over\tau_-^2}.
\end{equation}
For example $|\partial_c f|=d(\varkappa+v)/A^2$ and
$|\partial_\tau f|=d(\varkappa+v)^2/A^2$; the bounds follow from
$A\ge\tau(\varkappa+v)$ and $\varkappa+v\ge2a$.
The denominator stays positive throughout that algebraic parameter
rectangle.  A mean value estimate there interpolates numbers in
an explicit rational function, not physical trajectories.

Writing $\delta v_i=|v_i-\widehat v_i|$, this proves
\[
 \delta v_{i+1}\le\vartheta_s\delta v_i
   +{ |c_i-\widehat c_i|\over2a\tau_-^2}
   +{ |c_{i+1}-\widehat c_{i+1}|\over\tau_-}
   +{ |\tau_i-\widehat\tau_i|\over\tau_-^2}.
\]
The initial slope difference is at most $K-a$.  Sum the recurrence,
or sum its finite geometric expansion, and use
\eqref{eq:xd-summed-data}.  All terms are nonnegative and the initial
term is paid once.  Thus
\begin{equation}\label{eq:xd-slope-sum}
                    \sum_{i=0}^{m-1}\delta v_i\le D_v.
\end{equation}

The exact one-flight Euclidean Jacobian along the transported vector
is
\[
 {A_i\over c_{i+1}}
       {\sqrt{1+v_{i+1}^2}\over\sqrt{1+v_i^2}}.
\]
Multiplication therefore gives the actual identity
\begin{equation}\label{eq:xd-log-jacobian}
 \log\mathcal J_m(x,\xi)
 =\sum_{i=0}^{m-1}\log A_i-\sum_{i=1}^{m}\log c_i
          +{1\over2}\log{1+v_m^2\over1+v_0^2}.
\end{equation}
There is no incidence or Euclidean norm factor suppressed at the
two ends.  In particular multiplication by $c_m$ cancels just the
terminal logarithm in the second sum.

For the paired scalar numerators,
\[
 |A_i-\widehat A_i|
 \le(\varkappa_++K)|\tau_i-\widehat\tau_i|
      +\tau_+\delta v_i+|c_i-\widehat c_i|.
\]
Since both numerators are at least $2a\tau_-$,
\eqref{eq:xd-summed-data} and \eqref{eq:xd-slope-sum} imply
\begin{equation}\label{eq:xd-numerator-sum}
          \sum_{i=0}^{m-1}|\log A_i-\log\widehat A_i|\le D_A.
\end{equation}
After cancellation of the terminal incidence, the logarithmic
incidences are precisely the indices $1,\ldots,m-1$ covered by the
assumption.  Their total difference is at most $D_{\log}$ by
Lemma~\ref{lem:xd-log-incidence}.
The last term of \eqref{eq:xd-log-jacobian} telescopes before any
bound is applied.  The difference of its paired endpoint values
is at most
$\log((1+K^2)/(1+a^2))=D_{\rm end}$.
These three estimates prove \eqref{eq:xd-weighted-distortion} and
its exact equivalent \eqref{eq:xd-endpoint-residual}.

If the terminal index also has a common strip, apply
\eqref{eq:xd-log-sum} to all indices $1,\ldots,m$ instead.
Its constant was defined using the sum over $0,\ldots,m$, so no
additional endpoint constant is required.  This proves
\eqref{eq:xd-unweighted-distortion}.  For $m=1$ the interior sum is
empty, the slope-sum bound includes only the initial slope, and the
same proof gives both stated conclusions.  Only actual finite words
and their true reflection derivatives have been used.
\end{proof}

\begin{lemma}[A physical terminal-incidence limit]
\label{lem:xd-physical-terminal-limit}
In the original triangular lattice at $R=9/20$, there is a family
of regular one-flight records with the fixed lifted word
$(O_0,O_1)=((0,0),(1,0))$ for which $c_0$ stays strictly positive,
$c_1\downarrow0$, and, for vectors $\xi$ of fixed slope $a$,
\begin{equation}\label{eq:xd-physical-terminal-limit}
 c_1\mathcal J_1(x,\xi)\longrightarrow L_1\in(0,\infty),
 \qquad \mathcal J_1(x,\xi)\longrightarrow\infty.
\end{equation}
Every open flight segment in the family avoids every other obstacle.
Thus the endpoint factor retained in
\eqref{eq:xd-endpoint-residual} is present on actual trajectories,
including when the complete one-flight word is fixed.
\end{lemma}

\begin{proof}
Start at $Q_0=(R,0)$ with the outward unit velocity
$v_\theta=(\cos\theta,\sin\theta)$, where
\[
 0\le\theta<\theta_*:=\arcsin(9/11).
\]
The distance from $Q_0$ to $O_1$ is $11/20$.  Direct intersection
of this ray with the radius-$R$ circle gives its first positive
entry time and terminal incidence:
\begin{equation}\label{eq:xd-tangent-family}
 c_1(\theta)=\sqrt{1-(11\sin\theta/9)^2},\qquad
 \tau(\theta)={11\over20}\cos\theta-{9\over20}c_1(\theta).
\end{equation}
The time is positive: the difference between the squares of its
two nonnegative terms is
$(11/20)^2-(9/20)^2=1/10$.
The initial incidence is $c_0=\cos\theta>0$.

We check the entire physical segment, not just the entry point.
For $0\le t\le\tau(\theta)$, the bound
$\tau\le(11/20)\cos\theta$ puts $Q_0+tv_\theta$ in
\[
          [9/20,1]\mathbin\times[0,11/40].
\]
The nonzero rows of the triangular lattice have vertical coordinates
of absolute value at least $\sqrt3/2$.  Their distance from this
rectangle is therefore at least
$\sqrt3/2-11/40>9/20$; the strict inequality follows by squaring
the positive quantities $\sqrt3/2$ and $29/40$, whose squared
difference is $359/1600$.  On the zero row, every center other than
$(0,0),(1,0)$ is an integer horizontal translate with distance at
least $1$ from the rectangle.  The initial obstacle is avoided
for every $t>0$ because
\[
       |Q_0+tv_\theta|^2=R^2+2Rt\cos\theta+t^2>R^2.
\]
The smaller positive root in \eqref{eq:xd-tangent-family} is the
first contact with the target disk.  These facts exclude every
unvisited obstacle and every earlier contact.  For $\theta<\theta_*$
the flight is regular and has exactly the stated physical word.

As $\theta\uparrow\theta_*$,
\[
 c_0\longrightarrow {2\sqrt{10}\over11},\qquad
 \tau\longrightarrow {1\over\sqrt{10}},\qquad
 Q_1\longrightarrow
       \left({139\over220},{9\sqrt{10}\over110}\right).
\]
At each of these regular initial states take a nonzero vector of
slope $a$; a local smooth curve with that tangent lies in the same
regular one-flight domain after restriction to a sufficiently small
neighborhood.  Formula \eqref{eq:xd-A-and-f}, with
$\varkappa=20/9$, gives
\[
 A_0=\tau(\varkappa+a)+c_0,
 \qquad v_1=\varkappa+{c_1(\varkappa+a)\over A_0},
 \qquad c_1\mathcal J_1
      =A_0{\sqrt{1+v_1^2}\over\sqrt{1+a^2}}.
\]
Consequently \eqref{eq:xd-physical-terminal-limit} holds with
\[
 L_1=\left({\varkappa+a\over\sqrt{10}}
                    +{2\sqrt{10}\over11}\right)
                     {\sqrt{1+\varkappa^2}\over\sqrt{1+a^2}}>0.
\]
No derivative at the limiting grazing state is used.  Comparing this
family with its regular member $\theta=0$ shows directly why a
uniform unweighted comparison without terminal homogeneity cannot
be obtained by deleting the displayed terminal factor.  This is a
local geometric explanation of that factor; it is not an obstruction
to the whole-pressure sum or to the original raw local-limit target.
\end{proof}

\subsection{The precise pressure and optical consequences}

\begin{corollary}[The required radius-uniform transverse curve comparison]
\label{cor:xd-backward-curves}
Let $m\ge1$ and let $U,V$ be any two local smooth curves whose
tangent slopes at $u\in U$ and $v\in V$ belong to $[-K,-a]$.
Suppose the two actual backward records through times
$0,-1,\ldots,-m$ are regular, have the same lifted physical word,
and have matching homogeneity strips at times $-1,\ldots,-m$.
Then
\begin{equation}\label{eq:xd-backward-curves}
 C_\times^{-1}\le
 {J_UT_R^{-m}(u)\over J_VT_R^{-m}(v)}\le C_\times.
\end{equation}
In particular this holds on every one of the actual backward
homogeneous cells specified in the preserved pressure candidate,
with its literal time-index convention.  It does not require the
cell to meet a chosen transverse test.  If homogeneity is specified
only at $-1,\ldots,-m+1$, the exact endpoint-residual version is
\[
 C_\times^{-1}{c(T_R^{-m}v)\over c(T_R^{-m}u)}
 \le {J_UT_R^{-m}(u)\over J_VT_R^{-m}(v)}
 \le C_\times{c(T_R^{-m}v)\over c(T_R^{-m}u)}.
\]
\end{corollary}

\begin{proof}
The actual time reversal $I(s,\varphi)=(s,-\varphi)$ is a
Euclidean arclength isometry and satisfies $T_R^{-1}=IT_RI$ on
the regular domain.  It turns both stable-cone directions into
unstable-cone directions and the two backward records into two
forward physical records with their order of centers reversed.
The reflected angles lie in the closure of the same reflected strip,
so the boundary-valid estimate of Lemma~\ref{lem:xd-log-incidence}
applies.  This argument does not assume that the original disjoint
assignment of strip boundaries commutes with reflection.  Apply the
appropriate form of Theorem~\ref{thm:xd-cross-distortion}.
Norms before and after $I$
are identical.  The comparison includes every admissible curve
tangent and needs neither common curvature nor a connection between
the two curves.
\end{proof}

\begin{corollary}[Whole-word normalized stable and optical cocycles]
\label{cor:xd-normalized-cocycles}
Let $x,\widehat x\in M_R'$ have the same actual $m$-flight lifted
word and matching interior homogeneity strips $1,\ldots,m-1$.
No endpoint homogeneity, momentum cap, or positive incidence cutoff
is imposed.  For the true stable Euclidean arclength cocycle,
\begin{equation}\label{eq:xd-stable-endpoint}
 C_\times^{-1}\le
 {j_m(x)/c_0(x)\over j_m(\widehat x)/c_0(\widehat x)}
                                                        \le C_\times.
\end{equation}
For the original optical multiplier $\omega$,
\begin{equation}\label{eq:xd-optical-endpoints}
 (30000^2C_\times)^{-1}\le
 {\omega(x)/(c_0(x)c_m(x))\over
       \omega(\widehat x)/(c_0(\widehat x)c_m(\widehat x))}
                                            \le30000^2C_\times.
\end{equation}
These statements retain both physical endpoint factors explicitly.
In particular an endpoint incidence is not declared harmless merely
because a word or roof label is held fixed.
\end{corollary}

\begin{proof}
Start the backward comparison at $T_R^mx,T_R^m\widehat x$ and use
the true stable tangent directions there.  The backward Jacobians
are exactly $1/j_m(x),1/j_m(\widehat x)$, by invariance of the
stable line.  The endpoint-weighted comparison of
Theorem~\ref{thm:xd-cross-distortion}, after time reversal, compares
$c_0(x)/j_m(x)$ to $c_0(\widehat x)/j_m(\widehat x)$.
Taking reciprocals proves \eqref{eq:xd-stable-endpoint}.
The already proved optical comparison of module 313 is
\[
 {c_mj_m\over30000}\le\omega\le30000c_mj_m.
\]
Divide by the two actual endpoint incidences and compare the two
points.  One upper and one lower optical inequality cost $30000^2$,
not $30000$ and not a factor raised to the number of collisions.
This gives \eqref{eq:xd-optical-endpoints}.
\end{proof}

\begin{remark}[Completed geometric input and the remaining positive sum]
\label{rem:xd-pressure-scope}
Corollary~\ref{cor:xd-backward-curves} completes the radius-uniform
cross-curve Jacobian comparison that remained conditional in the
preserved whole-pressure candidate.  Its proof is native and uses
no radius-family interpretation of a fixed-table external theorem.
It uses the established global minimum as a deterministic comparison
only, and it includes all actual words, all positive incidences,
and the exact homogeneity indices needed by its Jacobian formula.

The actual whole-cylinder pressure
\[
 \sup_{R,m}\sum_{w\in\mathcal W_m(R)}
                  \sup_{x\in[w]_R\cap M_R'}j_m(x)
\]
is not bounded merely by comparing two points in one homogeneous
component.  The positive sum still needs the complete physical
barrier/continuation and planar refinement argument, or another
proved payment for every cell supremum.  In particular a component
with arbitrarily small measure is not paid by module 315's bulk
integrability theorem alone.  The long-cell assignment and short-cell
renewal in the retained candidate become valid conclusions only
after their remaining actual planar source hypotheses are proved.

No original section, return, displacement, collision or physical roof
test has been modified.  Any such test, and any Borel selector taking
values in $[0,1]$, may restrict the physical points compared here;
they are not differentiated.  This theorem itself is not a roof
density estimate, a trace theorem, a uniform twisted spectral bound,
or a proof of the original collision-normalized shortage limit.
It changes neither the $m^2$ normalization nor the prescribed order
of the original outer $B$ limit and collision limsup.
\end{remark}

